% Point to the main file
\documentclass[../../Main.tex]{subfiles}

% ========== Start the document ==========
\begin{document}

\section{Section 1.5 - Nested Quantifiers}

Let $\mathbb{P}$ be the set of prime integers

$$\mathbb{P} = \{ 2, 3, 5, 7, 11, \ldots \}$$

With this in mind, we can write the following in two ways:

\begin{center}
    Consider the statement "for all positive integers $n$, $6n - 1$ or $6n + 1$ is prime"
\end{center}

\begin{enumerate}
    \item $\forall n \in \mathbb{Z}^+, 6n - 1 \in \mathbb{P} \lor 6n + 1 \in \mathbb{P}$
    \item $\forall n \in \mathbb{Z}^+ \exists p \in \mathbb{P}, 6n - 1 = p \lor 6n + 1 = p$
\end{enumerate}

\begin{example}[Example of Nested Quantifiers]
    
    \begin{enumerate}
        \item
            For all real $x$ there exists a real number $y$ such that $xy = 1$

            $$\forall x \in \mathbb{R} \exists y \in \mathbb{R}, xy = 1$$

            \begin{center}
                \FALSE, since if $x = 0$, then $xy = 0 \not = 1$
            \end{center}

        \item 
            For all positive real $x$, there exists a real number $y$ such that $xy = 1$
            $$\forall x \in \mathbb{R}^+ \exists y \in \mathbb{R}, xy = 1$$
            \begin{center}
                \TRUE since if $x \in \mathbb{R}^+$, then $x > 0$, and if $y = \frac{1}{x}$, then $xy = x \frac{1}{x} = 1$
            \end{center}

            \begin{note}
                Given $x$, we can find $y$

                \yellowbox{$y$ depends on $x$}
            \end{note}

        But what happens if we switch $\forall$ and $\exists$?

        \item
            $\exists y \in \mathbb{R} \forall x \in \mathbb{R}^+, xy = 1$

            \begin{center}
                \FALSE since if $x = 1$, then $y = 1$, but if $x = 2$, then $y = \frac{1}{2}$, and $1 \not = \frac{1}{2}$
            \end{center}

            \begin{note}
                There is a single $y$ that works \textbf{FOR ALL} $x$

                \yellowbox{$y$ does not depend on $x$}
            \end{note}

        \item
            For all  positive integers $a$ and $b$, we have $ab > 0$

            $$\forall a \in \mathbb{Z}^+ \forall b \in \mathbb{Z}^+, ab >0$$

            \begin{note}
                This is the same as $\forall a, b \in \mathbb{Z}^+$
            \end{note}

            \begin{center}
                \TRUE since $a,b \in \mathbb{Z}^+$, then $a > 0$ and $b > 0$, hence $ab > 0$
            \end{center}

        \item
            There exists integers $a$, and $b$ such that $a^2 + b^2 = 464$
            $$\exists a \in \mathbb{Z} \exists b \in \mathbb{Z}, a^2 + b^2 = 464$$
            $$\exists b \in \mathbb{Z} \exists a \in \mathbb{Z}, a^2 + b^2 = 464$$

            \begin{center}
                \TRUE $a = 8$, $b = 20$, 
                
                \begin{align*}
                    a^2 + b^2 &= 8^2 + 20^2 \\
                    &= 64 + 400 \\
                    &= 464 \\
                \end{align*}
            
            \end{center}

        \begin{note}
            The order of the \textbf{SAME} quantifier doesn't matter, only for \textbf{DIFFERENT} quantifiers does it matter
        \end{note}

    \end{enumerate}

\end{example}

\begin{example}
    Rewrite the following using Quantifiers

    \begin{enumerate}
        \item 
            "There exists positive integers $a$ and $b$ such that $a^2 + b^2 = 15625$

            \begin{alphabet}
                \item 
                    Write as Quantifier

                    $$\exists a \in \mathbb{Z}^+ \exists b \in \mathbb{Z}^+, a^2 + b^2 = 15625$$

                \item
                    Write the negation

                    $$\forall a \in \mathbb{Z}^+ \forall b \in \mathbb{Z}^+, a^2 + b^2 \not = 15625$$

                    \begin{work}
                        \begin{align*}
                            a^2 + b^2 &= 5^6 \\
                            a^2 + b^2 &= 5^4 \cdot 5^2 \\
                        \end{align*}

                        \begin{align*}
                            3^2 + 4^2 &= 5^2 \\
                            5^4(3^2 + 4^2) &= 5^4 \cdot 5^2 \\
                            5^4 \cdot 3^2 + 5^4 \cdot 4^2 &= 5^6 \\
                            (5^2 \cdot 3)^2 + (5^2 \cdot 4)^2 &= 5^6 \\
                            (75)^2 + (100)^2 &= 5^6 \\
                        \end{align*}
                        
                        $$\therefore a = 75, b = 100 \\$$

                        \begin{center}
                            \TRUE
                        \end{center}

                    \end{work}

            \end{alphabet}

    \end{enumerate}
\end{example}

There's more examples but I'm just too lazy to write them
if you get the chance, they're on page 3 of your notes

% ========== End the document ==========
\end{document}